\documentclass{article}

\usepackage[utf8]{inputenc}
\usepackage[T1]{fontenc}
\usepackage{helvet}
\usepackage{geometry}
\usepackage{xcolor}
\usepackage{eso-pic}
\usepackage{fancyhdr}
\usepackage{parskip}
\usepackage{microtype}
\usepackage[english, ukrainian]{babel}
\usepackage{url}
\usepackage{amsmath}
\usepackage{amssymb}
\usepackage{amsthm}
\usepackage{longtable}
\usepackage{booktabs}
\usepackage{hyperref}
\usepackage{titlesec}
\usepackage{tocloft}

\renewcommand{\cftsecfont}{\small\bfseries}
\renewcommand{\cftsubsecfont}{\footnotesize}
\renewcommand{\cftsubsubsecfont}{\scriptsize}
\renewcommand{\cftsecpagefont}{\small\bfseries}
\renewcommand{\cftsubsecpagefont}{\footnotesize}
\renewcommand{\cftsubsubsecpagefont}{\scriptsize}

\definecolor{nato}{HTML}{293757}
\definecolor{footergray}{gray}{0.65}

\geometry{a4paper,left=3cm,right=3cm,top=3cm,bottom=3cm}
\renewcommand{\familydefault}{\sfdefault}
\setcounter{secnumdepth}{3}

\titleformat{\section}
  {\color{nato}\Huge\bfseries\raggedright}{\thesection.}{1em}{}
\titlespacing*{\section}{0pt}{30pt}{12pt}

\titleformat{\subsection}
  {\color{nato}\LARGE\bfseries\raggedright}{\thesubsection.}{1em}{}
\titlespacing*{\subsection}{0pt}{20pt}{8pt}

\titleformat{\subsubsection}
  {\color{nato}\Large\bfseries\raggedright}{\thesubsubsection.}{1em}{}
\titlespacing*{\subsubsection}{0pt}{10pt}{6pt}

\fancyhf{}
\fancyhead[R]{\small\color{footergray} 2 червня 2026}
\fancyfoot[L]{\small\color{footergray} Комплексна система захисту інформації. Галузевий профіль безпеки для VPN та PKI (20).}
\fancyfoot[R]{\small\color{footergray}\thepage}

\newcommand\HUGE[1]{\fontsize{40}{40}\selectfont #1}

\renewcommand{\headrulewidth}{0pt}
\renewcommand{\footrulewidth}{0.4pt}
\renewcommand{\footrule}{\color{footergray}\hrule width \textwidth height 0.4pt}

\begin{document}

\pagestyle{fancy}

\begin{titlepage}
  \AddToShipoutPictureBG*{\AtPageLowerLeft{\color{nato}\rule{\paperwidth}{\paperheight}}}
  \color{white}\thispagestyle{empty}\vspace*{4cm}\centering
  {\Huge  \textbf{Комплексна система захисту інформації} \par} \vspace{2cm}
  {\HUGE  \textbf{Галузевий профіль безпеки для VPN та PKI (20)} \par} \vspace{2.5cm}
  \vfill
  {\large 2026 \copyright\ Критпографічні Телесистеми \par}
\end{titlepage}

\newpage

\begin{center}
  {\Huge\color{nato}\bfseries Зміст\par}
\end{center}
\vspace{40pt}
\tableofcontents

\begin{abstract}
Галузевий профіль — розробляється для забезпечення безпеки VPN-продуктів та інфраструктури PKI (CA, OCSP, TSP, LDAP).
\end{abstract}

\section{AC}
Клас заходів захисту AC --- УПРАВЛІННЯ ДОСТУПОМ

Цей клас зосереджується на обмеженні доступу до інформаційних систем, ресурсів та функцій лише для авторизованих користувачів, програм і пристроїв.

\textbf{Перелік заходів захисту:} Управління обліковими записами (AC-2); Забезпечення доступу (AC-3); Id-spe-ac-17-1; Id-spe-ac-17-2; Id-spe-ac-17-9.

\subsection{УПРАВЛІННЯ ОБЛІКОВИМИ ЗАПИСАМИ (AC-2)}
a. Визначити та задокументувати типи облікових записів системи, дозволених для використання в ІС для підтримки цілей, завдань, функцій і процесів організації.\\ 
b. Призначити менеджерів облікових записів для управління системними обліковими записами.\\ 
c. Створити умови для групового та рольового членства.\\ 
d. Визначити авторизованих користувачів інформаційної системи, членство в групі та ролі, а також дозволи доступу (наприклад, привілеї) та інші атрибути (за потреби) для кожного облікового запису.\\ 
e. Вимагати схвалення [Призначення: визначеною організацією відповідальною особою або роллю] запитів на створення облікових записів системи.\\ 
f. Створювати, активувати, змінювати, деактивувати та видаляти системні облікові записи відповідно до [Призначення: визначених організацією політики, процедур та умов].\\ 
g. Впровадити моніторинг використання облікових записів системи.\\ 
h. Повідомляти адміністраторів облікових записів у межах [Призначення: визначеного організацією часового періоду для кожної ситуації]:\\ 
1. коли облікові записи більше не потрібні;\\ 
2. коли користувачі звільнені чи переведені;\\ 
3. коли використовуються індивідуальні системи або наявні зміни, які потребують нових знань.\\ 
i. Авторизувати доступ до системи на основі:\\ 
1. Дійсної авторизації доступу.\\ 
2. Передбачуваного використання системи.\\ 
3. Інших атрибутів, що вимагаються організацією.\\ 
j. Проводити перегляд облікових записів на відповідність вимогам управління обліковими записами з [Призначення: визначеною організацією частотою].\\ 
k. Впровадити процес повторного випуску облікових даних спільного/групового облікового запису (якщо він буде розгорнутий), коли особи виходять з групи.\\ 
l. Узгодити процеси управління обліковими записами з процесами звільнення та переводу (передачі повноважень) персоналу.

\vspace{10pt}
{\footnotesize\bfseries
No: 1\\
Name: ac\_2\_odp\_09\\
Type: integer\\
Default: 24
}

{\footnotesize Визначено передумови та критерії членства в групах і ролях; визначено атрибути (за необхідності) для кожного облікового запису; визначено персонал або ролі, необхідні для затвердження запитів на створення облікових записів; визначено політику, процедури, передумови та критерії створення, активації, зміни, деактивації та видалення облікових записів; визначено персонал або ролі, які мають бути повідомлені; визначено період часу, протягом якого адміністратори облікових записів повинні бути повідомлені про те, що облікові записи більше не потрібні; визначено термін, протягом якого необхідно повідомляти адміністраторів облікових записів про звільнення або переведення користувачів; визначено період часу, протягом якого необхідно повідомляти адміністраторів облікових записів про зміни у використанні системи або необхідність знати про зміни для окремої особи; визначено атрибути, необхідні для авторизації доступу до системи (за потреби); AC-02\_ODP[10] AC-02a.[01] AC-02a.[02] AC-02b AC-02с AC-02d.01 AC-02d.02 AC-02d.03[01] AC-02d.03[02] AC-02e AC-02f.[01] AC-02f.[02] AC-02f.[03] AC-02f.[04] AC-02f.[05] AC-02g AC-02h.01 AC-02h.02 AC-02h.03 AC-02i.01 AC-02i.02 AC-02i.03 визначено періодичність перегляду облікових записів; визначено та задокументовано типи облікових записів, дозволених для використання в системі; визначено та задокументовано типи облікових записів, які заборонено використовувати в системі; призначені менеджери облікових записів; необхідні умови та критерії для членства в групах та ролях; визначено авторизованих користувачів системи; вказано приналежність до групи або ролі; для кожного облікового запису вказуються повноваження доступу (тобто привілеї); атрибути (за необхідності) вказуються для кожного облікового запису; для запитів на створення облікових записів потрібні схвалення від персоналу або ролей; облікові записи створюються відповідно до політики, процедур, передумов та критеріїв; облікові записи активуються відповідно до політики, процедур, передумов та критеріїв; облікові записи змінюються відповідно до політики, процедур, передумов та критеріїв; облікові записи деактивуються відповідно до політики, процедур, передумов та критеріїв; облікові записи видаляються відповідно до політики, процедур, передумов та критеріїв; контролюється використання облікових записів; адміністратори облікових записів та персонал або ролі отримують повідомлення протягом періоду часу, коли облікові записи більше не потрібні; адміністратори облікових записів та персонал або ролі отримують повідомлення протягом періоду часу, коли користувачі звільнені чи переведені; адміністратори облікових записів та персонал або ролі отримують повідомлення протягом періоду часу, коли використовуються індивідуальні системи або наявні зміни, які потребують нових знань. доступ до системи здійснюється на підставі дійсної авторизації доступу; доступ до системи авторизується на основі передбачуваного використання системи; доступ до системи авторизовано на основі атрибутів (за необхідності); AC-02j AC-02k.[01] AC-02k.[02] AC-02l.[01] AC-02l.[02] облікові записи переглядаються на відповідність вимогам управління обліковими записами частота; створено процес повторного випуску облікових даних спільного доступу або групових облікових записів (якщо вони розгорнуті), коли користувачів вилучено з групи; впроваджено процес повторного випуску облікових даних спільного доступу або групових облікових записів (якщо вони розгорнуті), коли користувачів вилучено з групи; процеси управління обліковими записами узгоджуються з процесами звільнення персоналу; процеси управління обліковими записами узгоджуються з процесами переводу персоналу}




\subsection{ЗАБЕЗПЕЧЕННЯ ДОСТУПУ (AC-3)}
Застосовувати затверджені повноваження для логічного доступу до інформації та ресурсів системи відповідно до чинної політики (правил) управління доступом.

\vspace{10pt}
{\footnotesize\bfseries
No: 1\\
Name: ac\_3\_01\\
Type: list\\
Default: [\textquotedbl{}default\_deny\_rule\textquotedbl{}, \textquotedbl{}abac\_rule\_1\textquotedbl{}]
}

{\footnotesize Затверджені повноваження на логічний доступ до інформації та ресурсів системи виконуються відповідно до чинних політик(правил) управління доступом}




\subsubsection{id-spe-ac-17-1}


\vspace{10pt}
{\footnotesize\bfseries
No: 1\\
Name: ac\_17\_1\_01\\
Type: string\\
Default: nil
}

{\footnotesize Проводиться моніторинг методами віддаленого доступу}


{\footnotesize\bfseries
No: 2\\
Name: ac\_17\_1\_02\\
Type: string\\
Default: nil
}

{\footnotesize Проводиться управління методами віддаленого доступу}




\subsubsection{id-spe-ac-17-2}


\vspace{10pt}
{\footnotesize\bfseries
No: 1\\
Name: ac\_17\_2\_01\\
Type: string\\
Default: \textquotedbl{}AES-256-GCM\textquotedbl{}
}

{\footnotesize ВІДДАЛЕНИЙ ДОСТУП - ЗАХИСТ КОНФІДЕНЦІЙНОСТІ ТА ЦІЛІСНОСТІ ЗА ДОПОМОГОЮ ШИФРУВАННЯ МЕТА ОЦІНКИ: Визначте, чи інформаційна система криптографічні механізми}




\subsubsection{id-spe-ac-17-9}


\vspace{10pt}
{\footnotesize\bfseries
No: 1\\
Name: ac\_17\_9\_01\\
Type: integer\\
Default: 30
}

{\footnotesize Передбачена можливість відключення або деактивації віддаленого доступу до системи протягом періоду часу}


{\footnotesize\bfseries
No: 2\\
Name: ac\_17\_9\_odp\\
Type: integer\\
Default: 30
}

{\footnotesize Визначено період часу, протягом якого потрібно відключити або деактивувати віддалений доступ до системи}




\section{AU}
Клас заходів захисту AU --- АУДИТ ТА ПІДЗВІТНІСТЬ

Цей клас забезпечує можливість відстеження дій у системі, генерування, захист та аналіз записів аудиту для виявлення порушень політики безпеки.

\textbf{Перелік заходів захисту:} Id-spe-au-10-5.

\subsubsection{id-spe-au-10-5}


\vspace{10pt}
{\footnotesize\bfseries
No: 1\\
Name: au\_10\_5\_01\\
Type: string\\
Default: nil
}

{\footnotesize НЕСПРОСТОВНІСТЬ - ЦИФРОВІ ПІДПИСИ [Вилучено: Включено до SI-07]}




\section{IA}
Клас заходів захисту IA --- ІДЕНТИФІКАЦІЯ ТА АВТЕНТИФІКАЦІЯ

Цей клас відповідає за однозначне розпізнавання користувачів або пристроїв та підтвердження їхньої справжності перед наданням доступу до системи.

\textbf{Перелік заходів захисту:} Id-spe-ia-2-11; Id-spe-ia-2-12; Управління ідентифікацією (IA-4); Id-spe-ia-5-2; Id-spe-ia-7.

\subsubsection{id-spe-ia-2-11}


\vspace{10pt}
{\footnotesize\bfseries
No: 1\\
Name: ia\_2\_11\_01\\
Type: string\\
Default: nil
}

{\footnotesize ІДЕНТИФІКАЦІЯ ТА АВТЕНТИФІКАЦІЯ (КОРИСТУВАЧІВ ОРГАНІЗАЦІЇ) - ВІДДАЛЕНИЙ ДОСТУП - ОКРЕМИЙ ПРИСТРІЙ [Вилучено: Включено до ІА-02(06)]}




\subsubsection{id-spe-ia-2-12}


\vspace{10pt}
{\footnotesize\bfseries
No: 1\\
Name: ia\_2\_12\_01\\
Type: list\\
Default: [\textquotedbl{}admin\textquotedbl{}, \textquotedbl{}security\_officer\textquotedbl{}]
}

{\footnotesize Приймаються та електронним шляхом підтверджуються повноваження облікових даних особистої ідентифікації}




\subsection{УПРАВЛІННЯ ІДЕНТИФІКАЦІЄЮ (IA-4)}


\vspace{10pt}
{\footnotesize\bfseries
No: 1\\
Name: ia\_4\_a\\
Type: list\\
Default: [\textquotedbl{}admin\textquotedbl{}, \textquotedbl{}security\_officer\textquotedbl{}]
}

{\footnotesize Управління ідентифікаторами здійснюється шляхом отримання дозволу від персоналу або ролей на призначення ідентифікатора особі, групі, ролі або пристрою}


{\footnotesize\bfseries
No: 2\\
Name: ia\_4\_b\\
Type: list\\
Default: [\textquotedbl{}admin\textquotedbl{}, \textquotedbl{}security\_officer\textquotedbl{}]
}

{\footnotesize Управління ідентифікаторами здійснюється шляхом вибору ідентифікатора, який ідентифікує окрему особу, групу, ролі або пристрій}


{\footnotesize\bfseries
No: 3\\
Name: ia\_4\_c\\
Type: list\\
Default: [\textquotedbl{}admin\textquotedbl{}, \textquotedbl{}security\_officer\textquotedbl{}]
}

{\footnotesize Управління ідентифікаторами здійснюється шляхом призначення ідентифікатора особі, групі, ролі або пристрою; IA-04(d) ідентифікатори управляються шляхом запобігання повторному використанню ідентифікаторів впродовж період часу}


{\footnotesize\bfseries
No: 4\\
Name: ia\_4\_odp\_01\\
Type: list\\
Default: [\textquotedbl{}admin\textquotedbl{}, \textquotedbl{}security\_officer\textquotedbl{}]
}

{\footnotesize Визначено персонал або ролі, від яких необхідно отримати дозвіл на призначення ідентифікатора}


{\footnotesize\bfseries
No: 5\\
Name: ia\_4\_odp\_02\\
Type: integer\\
Default: 30
}

{\footnotesize Визначено період часу для запобігання повторному використанню ідентифікаторів}




\subsubsection{id-spe-ia-5-2}


\vspace{10pt}
\textit{Немає параметрів для цього контролю.}
\vspace{10pt}


\subsection{id-spe-ia-7}


\vspace{10pt}
{\footnotesize\bfseries
No: 1\\
Name: ia\_7\_01\\
Type: list\\
Default: [\textquotedbl{}default\_deny\_rule\textquotedbl{}, \textquotedbl{}abac\_rule\_1\textquotedbl{}]
}

{\footnotesize Впроваджено механізми автентифікації в криптографічний модуль, який відповідає вимогам чинних законів, виконавчих розпоряджень, директив, політик, правил, стандартів та рекомендацій для такої автентифікації}




\section{SC}
Клас заходів захисту SC --- ЗАХИСТ ІНФОРМАЦІЙНОЇ СИСТЕМИ ТА

Цей клас охоплює механізми захисту інформації під час її передачі мережами зв'язку, криптографічний захист та ізоляцію критичних компонентів.

\textbf{Перелік заходів захисту:} Id-spe-sc-7-7; Id-spe-sc-8-1; Відключення мережі (SC-10); Встановлення ключами (SC-12); Id-spe-sc-17; Автентифікація сесії (SC-23); Id-spe-sc-28-1; Id-spe-sc-49; Апаратний захист (SC-51).

\subsubsection{id-spe-sc-7-7}


\vspace{10pt}
{\footnotesize\bfseries
No: 1\\
Name: sc\_7\_7\_01\\
Type: string\\
Default: \textquotedbl{}автоматизований засіб моніторингу\textquotedbl{}
}

{\footnotesize Запобігається розділеному тунелюванню для віддалених пристроїв, що підключаються до систем організації, якщо розділене тунелюванню не захищено за допомогою засоби захисту}


{\footnotesize\bfseries
No: 2\\
Name: sc\_7\_7\_odp\\
Type: string\\
Default: nil
}

{\footnotesize Визначені гарантії безпечного прокладання розділеному тунелюванню}




\subsubsection{id-spe-sc-8-1}


\vspace{10pt}
{\footnotesize\bfseries
No: 1\\
Name: sc\_8\_1\_odp\\
Type: string\\
Default: nil
}

{\footnotesize Вибрано одне або декілька з наступних ЗНАЧЕНЬ ПАРАМЕТРІВ: \{запобігти несанкціонованому розголошенню інформації; виявити зміни в інформації\}}




\subsection{ВІДКЛЮЧЕННЯ МЕРЕЖІ (SC-10)}
Завершити з'єднання з мережею, яке пов'язане із сеансом зв'язку в кінці сеансу або після [Призначення: визначений організацією період часу] бездіяльності.

\vspace{10pt}
{\footnotesize\bfseries
No: 1\\
Name: sc\_10\_odp\\
Type: integer\\
Default: 30
}

{\footnotesize Визначено період бездіяльності, після якого система розриває мережеве з'єднання, пов'язане з сеансом зв'язку; SC08(05)\_ODP[02] мережеве з'єднання, пов'язане з сеансом зв'язку, розірвано в кінці сеансу або після період часу бездіяльності}




\subsection{ВСТАНОВЛЕННЯ КЛЮЧАМИ (SC-12)}
Встановити та управляти криптографічними ключами для криптографічних засобів, які використовуються в системі відповідно до [Призначення: визначені організацією вимоги до генерації, поширення, зберігання, доступу та знищення ключів].

\vspace{10pt}
{\footnotesize\bfseries
No: 1\\
Name: sc\_12\_01\\
Type: string\\
Default: \textquotedbl{}AES-256-GCM\textquotedbl{}
}

{\footnotesize Встановлюються криптографічні ключі, коли в системі використовується криптографія відповідно до < SC-12\_ODP вимог >}


{\footnotesize\bfseries
No: 2\\
Name: sc\_12\_02\\
Type: string\\
Default: \textquotedbl{}AES-256-GCM\textquotedbl{}
}

{\footnotesize Здійснюється управління криптографічними ключами, коли в системі використовується криптографія, відповідно до вимог }


{\footnotesize\bfseries
No: 3\\
Name: sc\_12\_odp\\
Type: string\\
Default: nil
}

{\footnotesize Визначені вимоги до генерації, розповсюдження, зберігання, доступу та знищення ключів}




\subsection{id-spe-sc-17}
a. Випускати сертифікати відкритого ключа відповідно до [Призначення: визначеної організацією політики сертифікації];\\ 
b. Отримувати сертифікати відкритого ключа від затвердженого постачальника послуг.

\vspace{10pt}
\textit{Немає параметрів для цього контролю.}
\vspace{10pt}


\subsection{АВТЕНТИФІКАЦІЯ СЕСІЇ (SC-23)}
Забезпечити автентифікацію сеансів зв'язку.

\vspace{10pt}
{\footnotesize\bfseries
No: 1\\
Name: sc\_23\_01\\
Type: string\\
Default: nil
}

{\footnotesize Захищено автентифікацію сеансів зв'язку}




\subsubsection{id-spe-sc-28-1}


\vspace{10pt}
{\footnotesize\bfseries
No: 1\\
Name: sc\_28\_1\_01\\
Type: string\\
Default: \textquotedbl{}AES-256-GCM\textquotedbl{}
}

{\footnotesize Реалізовані криптографічні механізми для запобігання несанкціонованому розкриттю інформації, що знаходиться в стані спокою на системних компонентах або носіях}


{\footnotesize\bfseries
No: 2\\
Name: sc\_28\_1\_02\\
Type: string\\
Default: \textquotedbl{}AES-256-GCM\textquotedbl{}
}

{\footnotesize Реалізовані криптографічні механізми для запобігання несанкціонованій модифікації інформації, що знаходиться в стані спокою на системних компонентах або носіях}




\subsection{id-spe-sc-49}
Впровадити механізми апаратного поділу та застосування політики між [Призначення: домени безпеки, визначені організацією].

\vspace{10pt}
{\footnotesize\bfseries
No: 1\\
Name: sc\_49\_01\\
Type: list\\
Default: [\textquotedbl{}default\_deny\_rule\textquotedbl{}, \textquotedbl{}abac\_rule\_1\textquotedbl{}]
}

{\footnotesize Впроваджено механізми апаратного розділення та застосування політик між доменами безпеки}


{\footnotesize\bfseries
No: 2\\
Name: sc\_49\_odp\\
Type: list\\
Default: [\textquotedbl{}default\_deny\_rule\textquotedbl{}, \textquotedbl{}abac\_rule\_1\textquotedbl{}]
}

{\footnotesize Визначені домени безпеки, які потребують апаратного розділення та механізмів забезпечення дотримання політики}




\subsection{АПАРАТНИЙ ЗАХИСТ (SC-51)}
a. Перевіряти правильність роботи [Призначення: визначені організацією функції безпеки та приватності].\\ 
b. Виконувати перевірку [Вибір (один або кілька): [Призначення: визначені організацією системні перехідні стани]; за командою користувача з відповідними повноваженнями; [Призначення: визначена організацією частота]].\\ 
c. Повідомляти [Призначення: визначені організацією персонал або посадові особи] про невдалі перевірки безпеки та приватності.\\ 
d. [Вибір (один або кілька): Вимкнути систему; Перезапустити систему; [Призначення: визначені організацією альтернативні дії]], коли виявляються аномалії.

\vspace{10pt}
{\footnotesize\bfseries
No: 1\\
Name: sc\_51\_odp\_01\\
Type: string\\
Default: nil
}

{\footnotesize Визначено компоненти системної прошивки, потребують апаратного захисту від запису; які}


{\footnotesize\bfseries
No: 2\\
Name: sc\_51\_odp\_02\\
Type: list\\
Default: [\textquotedbl{}admin\textquotedbl{}, \textquotedbl{}security\_officer\textquotedbl{}]
}

{\footnotesize Визначені уповноважені особи, які повинні виконувати процедури вимкнення та повторного ввімкнення апаратного захисту від запису; SC-51a. використовується апаратний захист від запису для компонентів мікропрограми системи; SC-51b.[01] впроваджено спеціальні процедури для уповноважених осіб для ручного вимкнення апаратного захисту від запису для модифікацій мікропрограми; SC-51b.[02] реалізовано спеціальні процедури для уповноважених осіб для повторного увімкнення захисту від запису перед поверненням до робочого режиму}




\end{document}
